\documentclass[12pt,a4paper,notitlepage]{article}
\usepackage{amsmath}
\usepackage{amssymb}
\usepackage{amsbsy}
\usepackage{mathtools}
\usepackage{float}
\usepackage[french]{babel}
\usepackage{tikz}

\usepackage{palatino}

 \usepackage[active]{srcltx}
\usepackage{scrtime}

\newcounter{qnumber}
\setcounter{qnumber}{0}

\newcounter{enumber}
\setcounter{enumber}{0}

\newcommand{\question}{\textbf{(\addtocounter{qnumber}{1}\theqnumber)}\,}
\newcommand{\exercice}{\textbf{\addtocounter{enumber}{1}\textsc{Exercice} \theenumber}.\,}
\setlength{\parindent}{0cm}

% Absorbe les petits dépassements de justification (lignes trop longues)
% en répartissant l'élasticité sur les espaces plutôt que dans la marge.
\setlength{\emergencystretch}{3em}


\begin{document}

\title{\textsc{Économétrie Approfondie}\\ \small{(Éléments de correction)}}
\date{Jeudi 18 décembre 2025}

\maketitle

\thispagestyle{empty}


\bigskip

\exercice On suppose que les données sont générées par le modèle suivant~:
\[
  y_i = x_{i,1}\beta_1 + \ldots + x_{i,K}\beta_K + \varepsilon_i
\]
où $x_{i,k}$, pour $k=1,\ldots,K$, sont des variables explicatives
déterministes, $\beta_k$, pour $k=1,\ldots,K$, sont des paramètres
réels, $\varepsilon_i$ une variable aléatoire i.i.d. centrée de
variance $\sigma_{\varepsilon}^2$.

\question On obtient la représentation matricielle de ce modèle en
concaténant (verticalement) les $N$ observations. On pose~:
$Y = \left( y_1, y_2, \dots, y_N \right)'$,
$\varepsilon = \left( \varepsilon_1, \varepsilon_2, \dots, \varepsilon_N \right)'$,
$\beta = \left( \beta_1, \beta_2, \dots, \beta_K \right)'$ et
\[
  X =
  \begin{pmatrix}
    x_{1,1} & x_{1,2} & \ldots & x_{1,K}\\
    x_{2,1} & x_{2,2} & \ldots & x_{2,K}\\
    \vdots  & \vdots  &        & \vdots \\
    x_{N,1} & x_{N,2} & \ldots & x_{N,K}\\
  \end{pmatrix}
\]
$X$ est une matrice $N\times K$ (chaque ligne rassemble les $K$
variables explicatives d'une observation, chaque colonne rassemble les
$N$ observations d'une variable explicative). On peut alors réécrire le
modèle sous la forme~:
\[
Y = X\beta + \varepsilon
\]

\question L'estimateur des MCO est le vecteur $\hat b$ qui minimise la
somme des carrés des résidus~:
\begin{equation*}
  \begin{split}
    \hat b &= \arg\min_{\{b\}} (Y-Xb)'(Y-Xb) \\
              &= \arg\min_{\{b\}} Y'Y - b'X'Y - Y'Xb + b'X'Xb \\
              &= \arg\min_{\{b\}}  b'X'Xb -2b'X'Y
  \end{split}
\end{equation*}
(le terme $Y'Y$ ne dépend pas de $b$ et les scalaires $b'X'Y$ et
$Y'Xb$ sont égaux). La condition nécessaire d'optimalité (en dérivant
par rapport au vecteur $b$) est~:
\[
2X'X\hat b - 2X'Y = 0
\]
\[
\Leftrightarrow \hat b = \left( X'X \right)^{-1}X'Y
\]
Pour que cet estimateur existe, il faut que la matrice $X'X$ soit
inversible. C'est le cas si, et seulement si, $X$ est de plein rang
colonne ($\mathrm{rg}(X)=K$), ce qui suppose en particulier qu'il n'y
a pas de colinéarité parfaite entre les variables explicatives et que
$N\geq K$. La condition du second ordre est satisfaite car $X'X$ est
semi-définie positive (définie positive sous l'hypothèse de plein
rang), de sorte que le point critique est bien un minimum.

\question $\hat b$ est un estimateur sans biais de $\beta$. Pour le
montrer, substituons le processus générateur des données (DGP) dans
l'expression de l'estimateur~:
\[
  \begin{split}
    \hat b &= \left( X'X \right)^{-1}X'\left( X\beta + \varepsilon \right)\\
               &= \left( X'X \right)^{-1}X'X\beta +\left( X'X \right)^{-1}X'\varepsilon\\
               &= \beta + \left( X'X \right)^{-1}X'\varepsilon
  \end{split}
\]
Comme les variables explicatives sont déterministes (sinon il faudrait
recourir au théorème des espérances itérées), on a donc~:
\[
  \begin{split}
    \mathbb E\left[\hat b\right] &= \beta + \mathbb E\left[\left( X'X \right)^{-1}X'\varepsilon\right]\\
                                    &= \beta + \left( X'X \right)^{-1}X'\mathbb E\left[\varepsilon\right]\\
                                    &= \beta
  \end{split}
\]
puisque $\mathbb E[\varepsilon]=0$.

\question La variance de $\hat b$ est définie par
$\mathbb V\left[\hat b\right] = \mathbb E\left[\left( \hat b-\mathbb E\left[ \hat b \right] \right)\left( \hat b-\mathbb E\left[ \hat b \right] \right)'\right]$. En
utilisant l'expression $\hat b-\beta=\left( X'X \right)^{-1}X'\varepsilon$
obtenue plus haut, et en notant que les $\varepsilon_i$ sont i.i.d.
(donc $\mathbb V[\varepsilon]=\sigma_\varepsilon^2 I_N$), il vient~:
\[
  \begin{split}
    \mathbb V\left[\hat b \right] &= \mathbb E\left[\left( X'X \right)^{-1}X'\varepsilon\varepsilon'X\left( X'X \right)^{-1}\right]\\
                                     &= \left( X'X \right)^{-1}X'\mathbb V\left[\varepsilon\right] X\left( X'X \right)^{-1} \\
                                     &= \left( X'X \right)^{-1}X'\sigma_{\varepsilon}^2 I_N X\left( X'X \right)^{-1} \\
                                     &= \sigma^2_{\varepsilon}\left( X'X \right)^{-1}X'X\left( X'X \right)^{-1}\\
                                     &= \sigma^2_{\varepsilon}\left( X'X \right)^{-1}
  \end{split}
\]
(les variables explicatives étant déterministes, la matrice
$\left( X'X \right)^{-1}X'$ sort de l'espérance).

\question L'estimateur des MCO est une variable aléatoire, et a donc
une variance, parce que $\hat b = \beta + \left( X'X \right)^{-1}X'\varepsilon$
est une fonction (linéaire) du vecteur aléatoire $\varepsilon$. La
matrice $X$ étant déterministe, c'est uniquement l'aléa des erreurs,
transmis par l'échantillon, qui rend $\hat b$ aléatoire. D'un
échantillon (d'une réalisation de $\varepsilon$) à l'autre, la valeur
de l'estimateur change~: cette dispersion autour de $\beta$ est
mesurée par la variance.

\question Pour montrer la convergence en probabilité, partons de~:
\[
\hat b = \beta + \left( X'X \right)^{-1}X'\varepsilon
       = \beta + \left( \frac{X'X}{N} \right)^{-1}\frac{X'\varepsilon}{N}
\]
On suppose que la matrice des moments empiriques des régresseurs
admet une limite finie et inversible~:
\[
\frac{X'X}{N} \xrightarrow[N\to\infty]{} Q
\]
où $Q$ est une matrice $K\times K$ définie positive. Par ailleurs, le
terme $\tfrac{1}{N}X'\varepsilon$ est d'espérance nulle (les $X$ sont
déterministes et $\mathbb E[\varepsilon]=0$), et sa variance tend vers
zéro~:
\[
\mathbb V\left[ \frac{X'\varepsilon}{N} \right] = \frac{\sigma_\varepsilon^2}{N}\,\frac{X'X}{N}
\xrightarrow[N\to\infty]{} 0
\]
Le vecteur $\tfrac{1}{N}X'\varepsilon$ converge donc en moyenne
quadratique, et par conséquent en probabilité, vers~$0$ (c'est une
application de la loi des grands nombres~: $\tfrac{1}{N}X'\varepsilon$
est une moyenne d'erreurs centrées). Par continuité de l'inversion
matricielle et du produit (théorème de Slutsky), il vient~:
\[
\hat b \xrightarrow[N\to\infty]{\mathrm{proba}} \beta + Q^{-1}\times 0 = \beta
\]
On peut aussi le voir directement~: la variance de l'estimateur
$\mathbb V[\hat b]=\sigma_\varepsilon^2\left( X'X \right)^{-1}
=\tfrac{\sigma_\varepsilon^2}{N}\left( X'X/N \right)^{-1}$ tend vers la
matrice nulle quand $N\to\infty$. Comme $\hat b$ est sans biais et que
sa variance s'annule à la limite, $\hat b$ converge en moyenne
quadratique, donc en probabilité, vers $\beta$~: l'estimateur des MCO
est convergent.

\setcounter{qnumber}{0}
\bigskip

\exercice On considère le modèle générateur des données~:
\[
y_i = x_i \beta + \varepsilon_i, \qquad i=1,\ldots,N,
\]
avec $\beta$ un vecteur colonne de paramètres et $x_i$ un vecteur ligne
de variables explicatives déterministes. Les erreurs sont normalement
distribuées~:
\[
\varepsilon = (\varepsilon_1, \ldots, \varepsilon_N)' \sim \mathcal{N}(\mathbf{0}, \Omega)
\]
avec $\Omega$ symétrique et définie positive. Sous forme matricielle,
$Y = X\beta + \varepsilon$, et le modèle empirique est $y_i = x_i b + \epsilon_i$.

\question Oui, l'estimateur des MCO $\hat b = \left( X'X \right)^{-1}X'Y$
est sans biais. En substituant le DGP~:
\[
\hat b = \beta + \left( X'X \right)^{-1}X'\varepsilon
\]
et comme $X$ est déterministe et $\mathbb E[\varepsilon]=0$~:
\[
\mathbb E\left[\hat b\right] = \beta + \left( X'X \right)^{-1}X'\mathbb E\left[\varepsilon\right] = \beta
\]
L'absence de biais ne dépend pas de la structure de $\Omega$~: elle
tient uniquement à l'exogénéité (ici la nature déterministe) des
régresseurs et à la nullité de l'espérance des erreurs.

\question En reprenant $\hat b - \beta = \left( X'X \right)^{-1}X'\varepsilon$
et en utilisant $\mathbb V[\varepsilon]=\Omega$~:
\[
  \begin{split}
    \mathbb V\left[\hat b \right] &= \left( X'X \right)^{-1}X'\,\mathbb V\left[\varepsilon\right]\, X\left( X'X \right)^{-1} \\
                                     &= \left( X'X \right)^{-1}X'\Omega X\left( X'X \right)^{-1}
  \end{split}
\]
Contrairement au cas sphérique ($\Omega=\sigma^2 I_N$), la matrice
$\Omega$ ne se simplifie plus avec les $X'X$~: on obtient une
expression «~en sandwich~».

\question Non, l'estimateur des MCO n'est pas efficace dans la classe
des estimateurs linéaires et sans biais. Le théorème de Gauss-Markov
ne s'applique en effet que lorsque la matrice de variance-covariance
des erreurs est scalaire ($\Omega=\sigma^2 I_N$, homoscédasticité et
absence d'autocorrélation). Ici $\Omega$ est quelconque (symétrique
définie positive)~: il existe un autre estimateur linéaire sans biais,
l'estimateur des moindres carrés généralisés (MCG, théorème
d'Aitken), dont la variance est plus faible (la différence des deux
matrices de variance est semi-définie positive). Les MCO ne sont donc
pas BLUE.

\question Comme $\Omega$ est symétrique définie positive, son inverse
l'est aussi et admet une factorisation $\Omega^{-1}=P'P$ (par exemple
de Cholesky), avec $P$ une matrice $N\times N$ inversible. En
prémultipliant le modèle par $P$~:
\[
\underbrace{PY}_{\tilde Y} = \underbrace{PX}_{\tilde X}\beta + \underbrace{P\varepsilon}_{\tilde\varepsilon}
\]
le modèle transformé a des erreurs sphériques~:
\[
\mathbb V[\tilde\varepsilon] = P\,\mathbb V[\varepsilon]\,P' = P\Omega P' = P (P'P)^{-1}P' = I_N
\]
On peut donc appliquer les MCO au modèle transformé~: l'estimateur
obtenu est alors efficace (théorème de Gauss-Markov). C'est
l'estimateur des MCG~:
\[
  \begin{split}
    \hat b_{\textsc{mcg}} &= \left( \tilde X'\tilde X \right)^{-1}\tilde X'\tilde Y\\
                          &= \left( X'P'PX \right)^{-1}X'P'PY\\
                          &= \left( X'\Omega^{-1}X \right)^{-1}X'\Omega^{-1}Y
  \end{split}
\]
Sa variance s'obtient comme celle d'un estimateur des MCO appliqué au
modèle transformé (dont les erreurs ont pour variance $I_N$)~:
\[
\mathbb V\left[\hat b_{\textsc{mcg}}\right] = \left( \tilde X'\tilde X \right)^{-1} = \left( X'\Omega^{-1}X \right)^{-1}
\]
On vérifie qu'il s'agit bien d'un estimateur sans biais~:
\[
\hat b_{\textsc{mcg}} = \beta + \left( X'\Omega^{-1}X \right)^{-1}X'\Omega^{-1}\varepsilon
\]
d'espérance $\beta$.

\question Si $\Omega$ est inconnue, on ne peut pas l'estimer
librement~: elle comporte $N(N+1)/2$ paramètres distincts, soit plus
que le nombre d'observations. Il faut donc en spécifier la structure,
c'est-à-dire faire dépendre $\Omega$ d'un petit nombre de paramètres~:
$\Omega=\Omega(\theta)$ (par exemple une forme d'hétéroscédasticité
$\sigma_i^2 = h(z_i,\theta)$, ou une structure d'autocorrélation des
erreurs). On procède alors en deux étapes (MCG \emph{réalisables},
FGLS)~:
\begin{itemize}
\item estimer le modèle par les MCO, récupérer les résidus
  $\hat\epsilon_i$ et en déduire une estimation $\hat\theta$ des
  paramètres de $\Omega$, d'où une matrice estimée $\hat\Omega=\Omega(\hat\theta)$~;
\item calculer l'estimateur des MCG réalisables en remplaçant $\Omega$
  par $\hat\Omega$~:
  \[
  \hat b_{\textsc{fgls}} = \left( X'\hat\Omega^{-1}X \right)^{-1}X'\hat\Omega^{-1}Y
  \]
\end{itemize}
Cet estimateur n'est en général plus sans biais en échantillon fini
(car $\hat\Omega$ dépend des données, donc de $\varepsilon$), mais il
reste convergent et asymptotiquement équivalent aux MCG si $\hat\Omega$
converge vers $\Omega$. On peut éventuellement itérer la procédure
jusqu'à convergence des estimations.

\question Oui, l'estimation par maximum de vraisemblance est possible
puisque la loi des erreurs est spécifiée (gaussienne). Comme
$Y\sim\mathcal N(X\beta,\Omega)$, la log-vraisemblance est~:
\[
\ln L(\beta,\Omega) = -\frac{N}{2}\ln 2\pi - \frac{1}{2}\ln|\Omega|
 - \frac{1}{2}\left( Y-X\beta \right)'\Omega^{-1}\left( Y-X\beta \right)
\]
À $\Omega$ donnée, maximiser $\ln L$ par rapport à $\beta$ revient à
minimiser la forme quadratique
$\left( Y-X\beta \right)'\Omega^{-1}\left( Y-X\beta \right)$. La
condition du premier ordre
\[
X'\Omega^{-1}\left( Y-X\beta \right)=0
\]
redonne exactement l'estimateur des MCG~:
\[
\hat\beta_{\textsc{mv}} = \left( X'\Omega^{-1}X \right)^{-1}X'\Omega^{-1}Y = \hat b_{\textsc{mcg}}
\]
Lorsque $\Omega$ est inconnue (mais paramétrée par $\theta$), on
maximise conjointement $\ln L$ par rapport à $\beta$ et $\theta$. Sous
normalité, le maximum de vraisemblance et les MCG coïncident donc pour
la partie $\beta$.

\setcounter{qnumber}{0}
\bigskip

\exercice Sur le marché d'un bien, la demande et l'offre à la date $t$
sont données par le modèle structurel~:
\[
  \begin{cases}
    y_t^d = \alpha_0 + \alpha_1 P_t + \alpha_2 R_t + u_t\\
    y_t^o = \gamma_0 + \gamma_1 P_t + \gamma_2 W_t + v_t
  \end{cases}
\]
où $P_t$ est le prix, $R_t$ le revenu des consommateurs, $W_t$ les
coûts de production, $u_t$ et $v_t$ des chocs i.i.d., centrés, de
variances $\sigma_u^2$ et $\sigma_v^2$ et non corrélés. Les échanges
ont lieu à l'équilibre ($y_t^d=y_t^o\equiv y_t$). Les coûts de
production sont orthogonaux au choc de demande
($\mathrm{cov}(W_t,u_t)=0$) et le revenu est orthogonal aux deux chocs.

\question Signes attendus. Pour la demande~: $\alpha_1<0$ (la demande
décroît avec le prix) et $\alpha_2>0$ (la demande croît avec le revenu,
bien normal). Pour l'offre~: $\gamma_1>0$ (l'offre croît avec le prix)
et $\gamma_2<0$ (une hausse des coûts de production réduit l'offre).
Les constantes $\alpha_0$ et $\gamma_0$ n'ont pas de signe attendu a
priori.

\question À l'équilibre, $y_t^d=y_t^o$, donc le prix doit être tel
que~:
\[
\alpha_0 + \alpha_1 P_t + \alpha_2 R_t + u_t = \gamma_0 + \gamma_1 P_t + \gamma_2 W_t + v_t
\]
En regroupant les termes en $P_t$~:
\[
(\alpha_1 - \gamma_1)P_t = (\gamma_0-\alpha_0) + \gamma_2 W_t - \alpha_2 R_t + v_t - u_t
\]
soit, comme $\alpha_1\neq\gamma_1$ (les pentes d'offre et de demande
sont de signes opposés)~:
\[
P_t = \frac{(\gamma_0-\alpha_0) + \gamma_2 W_t - \alpha_2 R_t + v_t - u_t}{\alpha_1 - \gamma_1}
\]
Le prix d'équilibre est une fonction (affine) des chocs structurels
$u_t$ et $v_t$~: c'est donc une variable aléatoire. C'est la
\emph{forme réduite} de l'équation du prix. (En substituant $P_t$ dans
l'une des équations, on obtiendrait de même la quantité d'équilibre
$y_t$, elle aussi aléatoire et corrélée aux deux chocs.)

\question Non, ce n'est pas une bonne idée. Le régresseur $P_t$ de
l'équation de demande est corrélé avec le terme d'erreur $u_t$ de
cette même équation. En effet, comme $u_t$ et $v_t$ ne sont pas
corrélés~:
\[
\mathrm{cov}(P_t,u_t) = \mathrm{cov}\!\left( \frac{v_t-u_t}{\alpha_1-\gamma_1}, u_t \right)
= \frac{-\sigma_u^2}{\alpha_1-\gamma_1} = \frac{\sigma_u^2}{\gamma_1-\alpha_1} \neq 0
\]
Le prix est donc une variable \emph{endogène}~: un choc de demande
$u_t$ déplace la courbe de demande, ce qui modifie le prix
d'équilibre. La condition d'orthogonalité entre régresseur et erreur
n'est pas vérifiée, et l'estimateur des MCO de l'équation de demande
est biaisé et non convergent (biais de simultanéité~/ d'endogénéité).
En revanche, le revenu $R_t$, supposé orthogonal aux chocs, n'est pas
problématique.

\question Il faut recourir à une méthode de variables instrumentales
(VI), estimée par les doubles moindres carrés (2SLS). Pour instrumenter
le prix $P_t$, il faut une variable corrélée avec $P_t$ mais
orthogonale au choc de demande $u_t$. Les coûts de production $W_t$
sont un instrument idéal~:
\begin{itemize}
\item $W_t$ est corrélé avec $P_t$, car $W_t$ apparaît dans la forme
  réduite du prix (avec le coefficient $\gamma_2/(\alpha_1-\gamma_1)\neq 0$,
  l'instrument est pertinent)~;
\item $W_t$ est orthogonal à $u_t$ par hypothèse (validité de
  l'instrument)~: les coûts de production n'affectent la demande qu'à
  travers le prix.
\end{itemize}
L'équation de demande $y_t = \alpha_0 + \alpha_1 P_t + \alpha_2 R_t + u_t$
contient une variable endogène ($P_t$) et une variable exogène incluse
($R_t$, qui s'instrumente elle-même). On dispose d'un instrument exclu
($W_t$)~: le modèle est \emph{juste identifié}. Concrètement, en notant
$Z_t=(1,R_t,W_t)$ le vecteur des instruments et $X_t=(1,P_t,R_t)$ les
régresseurs, l'estimateur des VI est~:
\[
\hat\alpha_{\textsc{vi}} = \left( Z'X \right)^{-1}Z'Y
\]
(ou, de façon équivalente, on régresse d'abord $P_t$ sur $(1,R_t,W_t)$
pour obtenir $\hat P_t$ — première étape — puis on régresse $y_t$ sur
$(1,\hat P_t,R_t)$ — seconde étape). Cet estimateur est convergent car
les instruments sont valides et pertinents~:
$\tfrac{1}{T}Z'u\xrightarrow{\mathrm{proba}}0$ et $\tfrac{1}{T}Z'X$
converge vers une matrice inversible.

\question Non, l'estimateur des variables instrumentales n'est pas sans
biais en échantillon fini. Il s'écrit comme un rapport de variables
aléatoires~:
\[
\hat\alpha_{\textsc{vi}} = \alpha + \left( Z'X \right)^{-1}Z'u
\]
et l'espérance de ce rapport n'est pas le rapport des espérances~: en
général,
\[
\mathbb E\big[(Z'X)^{-1}Z'u\big]\neq 0
\]
car $Z'X$ et $Z'u$ ne sont pas indépendants. L'estimateur des VI est seulement
\emph{convergent} (sans biais asymptotiquement)~: c'est précisément
parce qu'on accepte de renoncer à l'absence de biais en petit
échantillon que l'on peut corriger le problème d'endogénéité.

\bigskip

\textbf{\textsc{Illustration graphique.}} On représente les couples
$(y_t,P_t)$ effectivement observés (les équilibres de marché) dans le
plan quantité-prix. À chaque date, l'équilibre est l'intersection
d'une courbe de demande et d'une courbe d'offre, toutes deux déplacées
par leurs chocs respectifs.

\begin{figure}[H]
  \centering
  \begin{minipage}[t]{0.49\textwidth}
    \centering
    \begin{tikzpicture}[scale=0.78]
      \clip (-0.4,-0.4) rectangle (5.7,5.6);
      \draw[->] (0,0) -- (5.4,0) node[below right] {\footnotesize $y$};
      \draw[->] (0,0) -- (0,5.4) node[left] {\footnotesize $P$};
      % Courbes de demande (pente -1), déplacées par u
      \foreach \a in {4.0,4.8,5.6}{
        \draw[blue!55] (0,\a) -- (5,{\a-5});
      }
      % Courbes d'offre (pente +0.5), déplacées par v et W
      \foreach \b in {0.6,1.4,2.2}{
        \draw[red!55] (0,\b) -- (5,{\b+2.5});
      }
      % Nuage des équilibres
      \foreach \pt in {(2.267,1.733),(1.733,2.267),(1.2,2.8),%
        (2.8,2.0),(2.267,2.533),(1.733,3.067),%
        (3.333,2.267),(2.8,2.8),(2.267,3.333)}{
        \fill \pt circle (1.7pt);
      }
      % Droite des MCO : P = 3.1 - 0.25 y
      \draw[very thick] (0.5,2.975) -- (4.3,2.025);
      \node[blue!70] at (1.55,4.5) {\footnotesize demande};
      \node[red!70] at (4.1,4.55) {\footnotesize offre};
      \node at (3.95,1.75) {\footnotesize MCO};
    \end{tikzpicture}
    \\[2pt]
    {\footnotesize (a) Les deux courbes se déplacent.}
  \end{minipage}
  \hfill
  \begin{minipage}[t]{0.49\textwidth}
    \centering
    \begin{tikzpicture}[scale=0.78]
      \clip (-0.4,-0.4) rectangle (5.7,5.6);
      \draw[->] (0,0) -- (5.4,0) node[below right] {\footnotesize $y$};
      \draw[->] (0,0) -- (0,5.4) node[left] {\footnotesize $P$};
      % Demande fixe : P = 5 - y
      \draw[blue,very thick] (0,5) -- (5,0);
      % Offre déplacée uniquement par W (demande stable)
      \foreach \b in {0.5,1.5,2.5}{
        \draw[red!55] (0,\b) -- (5,{\b+2.5});
      }
      % Équilibres : ils se situent sur la courbe de demande
      \foreach \pt in {(3.0,2.0),(2.333,2.667),(1.667,3.333)}{
        \fill \pt circle (1.7pt);
      }
      \node[blue!80] at (1.35,4.45) {\footnotesize demande};
      \node[red!70] at (4.15,4.55) {\footnotesize offre};
    \end{tikzpicture}
    \\[2pt]
    {\footnotesize (b) Seule l'offre se déplace (via $W$).}
  \end{minipage}
\end{figure}

Sur le panneau \textbf{(a)}, la demande est déplacée par les chocs
$u_t$ et l'offre par $v_t$ (et $W_t$)~: les équilibres forment un nuage
de points dont la dispersion ne reproduit ni la courbe de demande ni la
courbe d'offre. La droite des MCO de $y$ sur $P$ ajustée sur ce nuage a
une pente comprise entre celle de la demande et celle de l'offre (c'est
une moyenne des deux pentes structurelles, pondérée par les variances
des chocs)~: elle ne restitue donc pas la pente $\alpha_1$ de la
demande. C'est l'expression graphique du biais de simultanéité.

\medskip

Sur le panneau \textbf{(b)}, on isole la variation du prix provoquée par
l'instrument $W_t$~: lorsque seule l'offre se déplace (la demande
restant stable, $u_t$ orthogonal à $W_t$), les équilibres successifs
\emph{glissent le long de la courbe de demande}, qui devient alors
identifiable. C'est exactement ce que réalise l'estimateur des
variables instrumentales~: il n'exploite que la part des fluctuations
du prix imputable à $W_t$, qui est orthogonale au choc de demande, pour
estimer de façon convergente la pente de la demande.

\end{document}

%%% Local Variables:
%%% mode: latex
%%% TeX-master: t
%%% End:
